% !TeX root = ../../tesi.tex

\section{Ecosistema tecnologico}
\label{sec:ecosistema}

Le interfacce dei negozi Shopify rivolte ai clienti, dette \gloss{storefront}{\emph{storefront}}, si
possono sviluppare in due modi:

\begin{itemize}
        \item con un tema della piattaforma: Shopify genera le pagine a partire da modelli scritti in
	\gloss{liquid}{Liquid}, il suo linguaggio di \emph{templating}, e dai dati del negozio;
        \item con un’architettura \gloss{headless}{\emph{headless}}: l’interfaccia è un’applicazione
	separata, costruita per esempio con \gloss{hydrogen}{Hydrogen}, il \emph{framework} di Shopify
	basato su React, e accede al negozio tramite \gloss{api}{API}, cioè interfacce con cui un
	programma mette a disposizione di altri programmi dati e operazioni, secondo un contratto definito.
\end{itemize}

Shopify offre due famiglie di API: le Admin API, per gestire i dati e le operazioni del negozio, e
le Storefront API, per catalogo, carrello e acquisto. Entrambe usano \gloss{graphql}{GraphQL}, un
linguaggio di interrogazione in cui il \emph{client} chiede con una sola richiesta esattamente i
campi che gli servono. Un tema Liquid riceve invece i dati del negozio durante la generazione delle
pagine. Per aggiungere funzionalità, Halue sviluppa applicazioni personalizzate in Node.js e
TypeScript, integrabili con entrambe le modalità.

Nei progetti di intelligenza artificiale cambia soprattutto l’architettura. Al centro c’è un
\gloss{llm}{modello linguistico di grandi dimensioni (LLM)}, addestrato a elaborare e generare
testo, che chiede all’agente che lo ospita di invocare uno strumento, ne riceve il risultato e
ripete il ciclo finché non formula la risposta (\figref{fig:agente}). Gli strumenti si possono
scrivere su misura per il progetto oppure ottenere tramite \gloss{mcp}{MCP}, un protocollo aperto
con cui un \emph{server} espone strumenti e risorse, come il catalogo di un negozio, a qualunque
agente lo supporti, senza un’integrazione dedicata per ciascuno.

In ChatCart ho sviluppato una ricerca nel catalogo su misura al posto del \emph{server} MCP di
Shopify: da questa scelta nasce il contratto tra agente e \emph{widget} che descrivo nella
\secref{sec:contratto}.

\begin{figure}[htbp]
	\centering
	\setstretch{1}
	% L'architettura di un progetto di intelligenza artificiale in Halue: l'applicazione dell'agente
% sta tra l'utente, il modello e i servizi esterni da cui arrivano dati e azioni.
% In bordeaux le parti sviluppate per ogni progetto.
% Il disegno è largo quanto il testo (150 mm): si include con \input, senza ridimensionarlo.

\begin{tikzpicture}[
	every node/.style={text=FiguraTesto},
	blocco/.style={draw=FiguraTenue, fill=white, rounded corners=3pt, text width=#1-2mm,
		minimum width=#1, minimum height=11mm, align=center, inner sep=1.5pt,
		font=\fontsize{7.5}{9}\selectfont,
		execute at begin node={\hyphenpenalty=10000\exhyphenpenalty=10000}},
	blocco/.default=40mm,
	interno/.style={draw=FiguraAccento, line width=0.8pt},
	doppia/.style={freccia, {Stealth[length=1.8mm,width=1.4mm]}-{Stealth[length=1.8mm,width=1.4mm]}},
	etichetta/.style={font=\fontsize{7.5}{9}\selectfont, text=FiguraTenue, inner sep=1pt,
		align=center},
]
	% L'applicazione dell'agente e i servizi esterni, come contenitori sullo sfondo.
	\begin{scope}[on background layer]
		\path[partizione] (32mm, -24mm) rectangle (98mm, -76mm);
		\path[intestazione] (32mm, -24mm) rectangle (98mm, -30mm);
		\node[titolo] at (65mm, -27mm) {Applicazione dell’agente};
		\path[partizione] (106mm, -24mm) rectangle (150mm, -76mm);
		\path[intestazione] (106mm, -24mm) rectangle (150mm, -30mm);
		\node[titolo] at (128mm, -27mm) {Servizi esterni};
	\end{scope}

	% Il modello, fuori dall'applicazione che lo usa.
	\node[blocco=66mm] (llm) at (65mm, -6.5mm)
		{\voce{Modello linguistico (LLM)}{servizio del fornitore, di solito Claude}};

	% Dentro l'applicazione.
	\node[blocco=54mm, interno] (istruzioni) at (65mm, -38.5mm)
		{\voce{Istruzioni}{compito, regole, tono del marchio}};
	\node[blocco=54mm, interno] (strumenti) at (65mm, -52.5mm)
		{\voce{Strumenti interni}{ricerche e azioni su misura del progetto}};
	\node[blocco=54mm] (client) at (65mm, -66.5mm)
		{\voce{\emph{Client} MCP}{scopre e invoca gli strumenti dei \emph{server}}};

	% Fuori dall'applicazione.
	\node[blocco=22mm, minimum height=16mm] (utente) at (11mm, -52.5mm)
		{\voce{Utente}{scrive in linguaggio naturale}};
	\node[blocco=40mm] (ucp) at (128mm, -38.5mm)
		{\voce{Negozio con UCP}{carrello, \emph{checkout}, pagamento}};
	\node[blocco=40mm] (api) at (128mm, -52.5mm)
		{\voce{API della piattaforma}{Shopify, Salesforce, gestionali}};
	\node[blocco=40mm] (mcp) at (128mm, -66.5mm)
		{\voce{\emph{Server} MCP}{della piattaforma o di terze parti}};

	% Il ciclo tra applicazione e modello.
	\draw[freccia] (58mm, -24mm) -- (58mm, -12mm);
	\draw[freccia] (72mm, -12mm) -- (72mm, -24mm);
	\node[etichetta, anchor=east] at (56.5mm, -18mm) {conversazione, istruzioni,\\descrizione degli strumenti};
	\node[etichetta, anchor=west] at (73.5mm, -18mm) {testo, oppure\\strumento da invocare};

	% Collegamenti con l'esterno.
	\draw[doppia] (utente.east) -- (32mm, -52.5mm);
	\draw[freccia] (strumenti.east) -- (api.west);
	\draw[freccia] (102mm, -52.5mm) |- (ucp.west);
	\fill[FiguraTenue] (102mm, -52.5mm) circle (0.5mm);
	\draw[freccia] (client.east) -- (mcp.west);
\end{tikzpicture}

	\caption{Esempio di architettura di un progetto di intelligenza artificiale.}
	\label{fig:agente}
\end{figure}
